\pagebreak
\section*{两个前馈层 = 对参数的注意力}\label{sec:parameter_attention}

除了注意力层外，我们的模型还包含逐位置前馈网络（第 \ref{sec:ffn} 节），由两个线性变换组成，中间有一个 ReLU 激活。实际上，这些网络也可以看作是注意力的一种形式。将这样一个网络的公式与简单的点积注意力层的公式进行比较（省略偏置和缩放因子）：

\begin{align*}
    FFN(x, W_1, W_2) = ReLU(xW_1)W_2 \\
    A(q, K, V) = Softmax(qK^T)V
\end{align*}

基于这些公式的相似性，两层前馈网络可以看作是一种注意力，其中 key 和 value 是可训练参数矩阵 $W_1$ 和 $W_2$ 的行，且我们在兼容性函数中使用 ReLU 而不是 Softmax。

%the compatablity function is $compat(q, k_i) = ReLU(q \cdot k_i)$ instead of $Softmax(qK_T)_i$.

基于这种相似性，我们尝试用类似于模型其他地方使用的注意力层来替换逐位置前馈网络。多头参数注意力子层与第 \ref{sec:multihead} 节所述的多头注意力相同，除了每个注意力头的``key''和``value''输入是可训练的模型参数，而不是前一层的线性投影。这些参数按 $\sqrt{d_{model}}$ 的因子进行缩放，以更接近激活。

在我们的第一个实验中，我们用一个多头参数注意力子层替换了每个逐位置前馈网络，该子层有 $h_p=8$ 个头、key 维度 $d_{pk}=64$ 和 value 维度 $d_{pv}=64$，每个注意力头使用 $n_p=1536$ 个 key-value 对。该子层总共有 $2097152$ 个参数，包括 query 投影和输出投影中的参数。这与我们替换的逐位置前馈网络中的参数数量相匹配。虽然理论计算量相同，但实际上注意力版本的步骤时间增加了约 30\%。

在我们的第二个实验中，我们使用了 $h_p=8$ 个头和每个注意力头 $n_p=512$ 个 key-value 对，同样与基础模型中的参数总数相匹配。

第一个实验的结果略差于基础模型，第二个实验的结果略好，见表 \ref{tab:parameter_attention}。

\begin{table}[h]
\caption{用多头参数注意力替换逐位置前馈网络会产生与基础模型类似的结果。所有指标都在英-德翻译开发集 newstest2013 上。}
\label{tab:parameter_attention}
\begin{center}
\vspace{-2mm}
%\scalebox{1.0}{
\begin{tabular}{c|cccccc|cccc}
\hline\rule{0pt}{2.0ex}
 & \multirow{2}{*}{$\dmodel$} & \multirow{2}{*}{$\dff$} &
\multirow{2}{*}{$h_p$} & \multirow{2}{*}{$d_{pk}$} & \multirow{2}{*}{$d_{pv}$} &
 \multirow{2}{*}{$n_p$} &
 PPL & BLEU & params & training\\
 & & & & & &  & (dev) & (dev) & $\times10^6$ & time \\
\hline\rule{0pt}{2.0ex}
base & 512 & 2048 & & & & & 4.92 & 25.8 & 65 & 12 hours\\
\hline\rule{0pt}{2.0ex}
AOP$_1$ & 512 & & 8 & 64 & 64 & 1536 & 4.92& 25.5  & 65 & 16 hours\\
AOP$_2$ & 512 & & 16 & 64 & 64 & 512 & \textbf{4.86} & \textbf{25.9}  & 65 & 16 hours \\
\hline
\end{tabular}
%}
\end{center}
\end{table}
